\subsection{Hom-categories}
\label{subsec:laxloc-homs}

\gutentag{00F7}%
For $X$, $Y \in \C$, the equivalence \eqref{eq:laxloc-adjunction} identifies
$R\colon \Hom_{\C}(X,Y) \to \Hom_{\D}(RX,RY)$
with
$\varepsilon_{X}^{\natural}\colon \Hom_{\C}(X,Y) \to \Hom_{\C}(LRX,Y)$,
naturally in $Y$, as
$R(k\varepsilon_{X}) \circ \eta_{RX} \simeq Rk$
by $\theta_{X}$. So $R$ is locally fully faithful if and only if every
$\varepsilon_{X}^{\natural}$
is fully faithful.

\begin{prop}
\label{prop:laxloc-homs}
\gutentag{00F8}%
Let $X \in \C$.
\begin{enumerate}
\item The transformation
$R\colon \Hom_{\C}(X,-) \Rightarrow \Hom_{\D}(RX,R-)$
of functors $\C \to \Cat$ has a right adjoint in $\func(\C,\Cat)$ if and only
if $\varepsilon_{X}$ has a left adjoint $l$, and the right adjoint is then
$l^{\natural}$, under
$\Hom_{\D}(RX,R-) \simeq \Hom_{\C}(LRX,-)$.
Its components are then coreflective if and only if
$l \dashv \varepsilon_{X}$
is a coreflection, e.g.\ if $R$ is locally fully faithful.
\item If $L \dashv R$ is a lax localization and $X = Ld$, then
$l = L\eta_{d}$
and
$l \dashv \varepsilon_{X}$
is a coreflection. That is, the functor
\begin{equation*}
\Hom_{\D}(d,RY) \longrightarrow \Hom_{\D}(Td,RY) \ , \qquad g \longmapsto
R\varepsilon_{Y} \circ Tg \ ,
\end{equation*}
is coreflective, in particular fully faithful, with right adjoint
$\eta_{d}^{\natural}$, naturally in $Y$.
\end{enumerate}
\end{prop}

\begin{proof}
\gutentag{00F9}%
(1) By \eqref{eq:laxloc-adjunction}, the transformation $R$ is
$\varepsilon_{X}^{\natural}$,
the image of $\varepsilon_{X}$ under the enriched
Yoneda embedding
$h\colon \C\opcat \to \func(\C,\Cat)$.
As in the proof of \cref{lem:laxloc-repr}, $h$ is fully faithful, takes
$l \dashv r$
to
$r^{\natural} \dashv l^{\natural}$
with the same unit and counit, and every adjunction between representables
arises this way; and a $2$-morphism of $\func(\C,\Cat)$ is invertible if and
only
if its components are. (2) By \cref{lem:laxloc-equiv}(4),
$L\eta_{d} \dashv \varepsilon_{Ld}$
is a coreflection, and under \eqref{eq:laxloc-adjunction},
$(L\eta_{d})^{\natural}$
becomes $\eta_{d}^{\natural}$ and $R$ becomes
$g \mapsto R(\varepsilon_{Y} \circ Lg) = R\varepsilon_{Y} \circ Tg$.
\end{proof}

\begin{exmp}
\label{ex:laxloc-homs}
\gutentag{00FA}%
Let $\D = \Cat$ and
$P = \{\emptyset \to \mathrm{pt}\}$,
so that
$\C := \Pc_{\emptyset,P}$
consists of the $\infty$-categories with an
initial object and the functors preserving it, and $R$ is the inclusion. As we
will see in \cref{ex:local-vs-pushout}(2), $R$ has the left adjoint
$L(d) = d^{\triangleleft}$,
and $L \dashv R$ is a normalized lax localization by \cref{cor:laxloc-P}.
\begin{enumerate}
\item For
$X = L(\mathrm{pt}) = \Delta^{1}$,
\cref{prop:laxloc-homs}(2) says
that the arrows $0 \to v$ out of an initial object form a coreflective full
subcategory of $\func(\Delta^{1},Y)$, with coreflector
$(u \to v) \mapsto (0 \to v)$.
For $Y = \Delta^{1}$ it is not reflective, as there is no map from
$(1 \to 1)$
to any
$(0 \to v)$.
\item Let $X$ be the walking retraction, with objects $0$, $1$ and morphisms
$i\colon 0 \to 1$,
$r\colon 1 \to 0$
with
$ri = \mathrm{id}_{0}$;
its object $0$ is initial. Then
$\varepsilon_{X}\colon X^{\triangleleft} \to X$
has no left adjoint $l$: for the cone point $\ast$,
$\mapp_{X}(1,\varepsilon_{X}(\ast)) = \mapp_{X}(1,0)$
is non-empty while only $\ast$ maps to $\ast$, so $l(1) = \ast$, and then $1$
would be initial. So
$R\colon \Hom_{\C}(X,-) \Rightarrow \func(X,-)$
has no right adjoint in $\func(\C,\Cat)$. Objectwise,
$\Hom_{\C}(X,Y) \simeq Y_{/0}$,
and $\func(X,Y)$ consists of the retractions
$(s,r)\colon A \rightleftarrows B$,
$rs \simeq \mathrm{id}_{A}$.
A morphism from
$(0 \rightleftarrows B',\ t\colon B' \to 0)$
to $(s,r)$ is a
$v\colon B' \to B$
with
$rv \simeq (0 \to A) \circ t$,
so $(s,r)$ has a coreflection if and only if the pullback of $r$ along
$0 \to A$
exists. Let $Y$ be the category generated by objects $0$, $A$, $B$, $C$ and
morphisms
$s\colon A \to B$,
$r\colon B \to A$,
$c\colon C \to 0$
and
$u,u'\colon C \to B$,
subject to $0$ being initial,
$rs = \mathrm{id}_{A}$
and
$ru = ru' = (0 \to A) \circ c$.
The morphisms to $0$ are $\mathrm{id}_{0}$ and $c$, the morphisms
$C \to B$ are $u$, $u'$ and
$(0 \to B) \circ c$,
and the endomorphisms of $C$ are $\mathrm{id}_{C}$ and
$(0 \to C) \circ c$.
So neither
$(0,\ 0 \to B)$
nor any $(C,v)$ is a pullback, and $(s,r)$ has no coreflection. It has no
reflection either, as there is no morphism $A \to 0$. So
$R\colon \Hom_{\C}(X,Y) \to \Hom_{\D}(RX,RY)$
is neither reflective nor coreflective.
\end{enumerate}
\end{exmp}

\begin{rmk}
\label{rmk:laxloc-continuous}
\gutentag{00FB}%
Following Johnstone--Joyal \cite{johnstonejoyal1982continuous}, who treat the
free completion under filtered colimits, call $X \in \C$ \emph{continuous} if
$\varepsilon_{X}$ has a left adjoint; for the ideal completion of posets, these
are the continuous dcpos. So \cref{prop:laxloc-homs}(1) says that
$R\colon \Hom_{\C}(X,-) \Rightarrow \Hom_{\D}(RX,R-)$
has a right adjoint if and only if $X$ is continuous. By
\cref{lem:laxloc-equiv}(4), every $Ld$ is continuous, and if $R$ is locally
fully faithful, every continuous $X$ is a retract of $LRX$, as
$l \dashv \varepsilon_{X}$
is then a coreflection. In \cref{ex:laxloc-homs}, $X$ is continuous if and only
if every object with a morphism to $0$ is initial, i.e.\ if and only if
$X \simeq Ld$
for some $d$, as we will see in \cref{ex:local-vs-pushout}(2); and retracts of
the $Ld$ are again of this form. So there continuity characterizes the
essential image of $L$. We do not know whether retracts of the $Ld$ are
continuous in general.
\end{rmk}

\begin{rmk}[the colax case]
\label{rmk:laxloc-colax}
\gutentag{00FC}%
Let $L \dashv R$ be a colax localization. Applying the results of this section
to the lax localization
$L^{\mathrm{co}} \dashv R^{\mathrm{co}}$
(\cref{subsec:laxloc-def}), and using that $(-)^{\mathrm{co}}$ exchanges
reflective and coreflective $1$-morphisms, and left Beck--Chevalley squares
with transposes of right Beck--Chevalley squares, gives the following.
We have
$\varepsilon L \dashv L\eta$
with counit $\varphi$ and
$\eta R \dashv R\varepsilon$
with unit $\theta^{-1}$, so $L$ makes the $\eta_{d}$ reflective and their
naturality squares $(g,Tg)$, drawn with $\eta_{d}$ horizontal, left
Beck--Chevalley; $\Sat_{L}$ is replaced by the $1$-morphisms that $L$ makes
reflective and the squares it makes left Beck--Chevalley. The modules form the
locally full sub-$(\infty,2)$-category $\D_{T}$ of the $x$ with $\eta_{x}$
coreflective and the $1$-morphisms whose transposed unit square is right
Beck--Chevalley (\cref{thm:kz-module-maps}), and
$\D_{T} = \Pc_{\Sigma_{\eta},\emptyset}$
in the notation that will be introduced in \cref{def:IP-local}:
the $i^{\natural}$
have fully faithful right adjoints. In (R1) and (R2), reflective $1$-morphisms
and their left adjoints are replaced by coreflective $1$-morphisms and their
right adjoints. The recognition criterion uses the lax limits
$(f \downarrow e)$
instead \cite[Rem.~3.3]{bourke2025recognition}. Finally,
$\Hom_{\C}(Ld,Y) \simeq \Hom_{\D}(d,RY)$
is a reflective full subcategory of
$\Hom_{\D}(Td,RY)$,
with reflector $\eta_{d}^{\natural}$.
\end{rmk}
